\documentclass[letterpaper,11pt]{article}

\usepackage[empty]{fullpage}
\usepackage{titlesec}
\usepackage[usenames,dvipsnames]{color}
\usepackage{enumitem}
\usepackage[T2A]{fontenc}
\usepackage[utf8]{inputenc}
\usepackage[hidelinks]{hyperref}
\usepackage{fancyhdr}
\usepackage[russian,english]{babel}
\usepackage{tabularx}
\input{glyphtounicode}

\pagestyle{fancy}
\fancyhf{} % Clear all header and footer fields
\fancyfoot{}
\renewcommand{\headrulewidth}{0pt}
\renewcommand{\footrulewidth}{0pt}

% Adjust margins
\addtolength{\oddsidemargin}{-0.5in}
\addtolength{\evensidemargin}{-0.5in}
\addtolength{\textwidth}{1in}
\addtolength{\topmargin}{-.5in}
\addtolength{\textheight}{1.0in}

\urlstyle{same}

\raggedbottom
\raggedright
\setlength{\tabcolsep}{0in}

% Sections formatting
\titleformat{\section}{
  \vspace{-4pt}\scshape\raggedright\large
}{}{0em}{}[\color{black}\titlerule \vspace{-5pt}]

% Ensure that generated PDF is machine readable/ATS parsable
\pdfgentounicode=1

%-------------------------
% Custom commands
\newcommand{\resumeItem}[2]{
  \item\small{
    \textbf{#1}{: #2 \vspace{-2pt}}
  }
}

% Detailed item: title + description, problem, result and tech stack each on its own line.
\newcommand{\resumeItemDetailed}[5]{
  \item\small{
    \textbf{#1}: #2 \\[3pt]
    \textit{Проблема}: #3 \\[3pt]
    \textit{Результат}: #4 \\[3pt]
    \textit{Стек}: #5
    \vspace{-2pt}
  }
}

% Summary item: title + what you did + result. Use for cross-cutting contributions
% like mentorship, processes, on-call, team culture -- where "problem"/"stack" don't apply.
\newcommand{\resumeSummary}[3]{
  \item\small{
    \textbf{#1}: #2 \\[3pt]
    \textit{Результат}: #3
    \vspace{-2pt}
  }
}

% Tech term: programming languages, libraries, tools, products.
% Wraps in monospace for visual emphasis and to make tech keywords scannable.
\newcommand{\tech}[1]{\texttt{#1}}

% Just in case someone needs a heading that does not need to be in a list
\newcommand{\resumeHeading}[4]{
    \begin{tabular*}{0.99\textwidth}[t]{l@{\extracolsep{\fill}}r}
      \textbf{#1} & #2 \\
      \textit{\small #3} & \textit{\small #4} \\
    \end{tabular*}\vspace{-5pt}
}

\newcommand{\resumeSubheading}[4]{
  \vspace{-1pt}\item
    \begin{tabular*}{0.97\textwidth}[t]{l@{\extracolsep{\fill}}r}
      \textbf{#1} & #2 \\
      \textit{\small #3} & \textit{\small #4} \\
    \end{tabular*}\vspace{-5pt}
}

\newcommand{\resumeSubSubheading}[2]{
    \begin{tabular*}{0.97\textwidth}{l@{\extracolsep{\fill}}r}
      \textit{\small #1} & \textit{\small #2} \\
    \end{tabular*}\vspace{-5pt}
}

\newcommand{\resumeItemPlain}[1]{
  \item\small{
    {#1 \vspace{-2pt}}
  }
}

\newcommand{\resumeSubItem}[2]{\resumeItem{#1}{#2}\vspace{-4pt}}

\renewcommand{\labelitemii}{$\circ$}

\newcommand{\resumeSubHeadingListStart}{\begin{itemize}[leftmargin=*]}
\newcommand{\resumeSubHeadingListEnd}{\end{itemize}}
\newcommand{\resumeItemListStart}{\begin{itemize}}
\newcommand{\resumeItemListEnd}{\end{itemize}\vspace{-5pt}}

%-------------------------------------------
%%%%%%  CV STARTS HERE  %%%%%%%%%%%%%%%%%%%%%%%%%%%%


\begin{document}

%----------HEADING-----------------
\begin{tabular*}{\textwidth}{l@{\extracolsep{\fill}}l}
  \textbf{\href{https://vanya.space/}{\Large Иван Ильин}} & Email: \href{mailto:buran05@icloud.com}{buran05@icloud.com}\\
  \href{https://vanya.space/}{vanya.space} & Mobile: \href{tel:+79033459111}{+7-903-345-9111} \\
   & Github: \href{https://github.com/avalanche05}{@avalanche05} \\
   & Telegram: \href{https://t.me/avalanche05}{@avalanche05} \\
\end{tabular*}

%--------SKILLS------------
\section{Навыки}
 \resumeSubHeadingListStart
   \item{
     \textbf{Языки}{: \tech{Python}, \tech{Golang}, \tech{C++}, \tech{Java}, \tech{SQL}}
     \hfill
     \textbf{Технологии}{: \tech{S3}, \tech{Kafka}, \tech{PostgreSQL}, \tech{ClickHouse}, \tech{Docker}, \tech{Kubernetes}}
   }
 \resumeSubHeadingListEnd

%-----------EXPERIENCE-----------------
\section{Опыт работы}
  \resumeSubHeadingListStart

    \resumeSubheading
      {Яндекс, Автономный Транспорт}{Москва, Россия}
      {Разработчик}{Сентябрь 2023 -- Сейчас}
      \resumeItemListStart
        \resumeItemDetailed{Система для записи данных с автономных машин}
          {Внедрил набор сервисов для записи данных с разных устройств в сети автономной машины.}
          {Внутри машины находится несколько компьютеров, которые суммарно генерируют поток данных ~500ГБ/час. Выгружать такой объём данных с машины занимает несколько часов.}
          {Съёмное устройство, которое позволяет записывать данные со всех устройств по сети. Это уменьшило простой машин с 4-х часов до 15-и минут в день.}
          {Клиент и сервер написаны на \tech{C++} с использованием Async Callback \tech{gRPC}. Сервисы для отправки метаданных в облака и реализации fallback сценариев написаны на \tech{Go}.}
        \resumeItemDetailed{Сервис для экстренного просмотра видео с автономной машины}
          {Разработал приложение для просмотра видео с машины в ситуациях, когда основной вычислитель машины повреждён.}
          {При повреждени основного компьютера на машине во время аварии, процесс получения видео занимает много времени и ручных действий}
          {Приложение, которое позволяет подключиться к хранилищу данных ноутбуком напрямую и прямо на месте инцидента получить видео с камер машины. Время получения видео сократилось с ~12 часов до 10-30 минут.}
          {Приложение читает \tech{rosbag}-и, в которых записан \tech{h265} поток видео, с помощью \tech{ffmpeg} это собирается в \tech{mp4} видео. Сервер преобразования видео написан на \tech{Go}. Клиент: \tech{html} templates + \tech{bootstrap} с ванильным \tech{js}.}
        \resumeItemDetailed{Obsevrability систем команды}
          {Настроил отправку метрик и алёрты.}
          {Тысячи устройств для записи данных, состояние работы каждого из них необходимо отслеживать и вовремя реагировать на проблемы.}
          {Дашборды, метрики и алёрты для отслеживания состояния работы устройств и сервисов на них.}
          {Все метрики отправляются в Яндексовый аналог \tech{Prometheus}, отрисовываются в \tech{DataLens}, алёрты настроены в Яндексовом аналоге \tech{Grafana}.}
        \resumeItemDetailed{Сервис для контроля данных на флоте автономных машин}
          {Расширил функциональность: добавил логику сохранения данных при смене локации для конкретного робота, улучшил UI для настройки конфигов.}
          {При смене локации робота-доставщика требовалось ручное вмешательство. Редактирование конфигов было неудобным, приходилось копировать конфиг в сторонние редакторы.}
          {Придумал и реализовал логику работы сервиса при смене локации робота-доставщика. Подключил плагин для корректного отображения \tech{yaml} конфигов. Запуск новых локаций стал независимым от сервиса, время редактирования конфигов сократилось с 10 минут до 1 минуты.}
          {Бэкенд управления написан на \tech{Python} \tech{Django}. Метаданные хранятся в \tech{PostgreSQL} с \tech{master/replica}-роутингом,
          аналитика событий в \tech{ClickHouse}.}
        \resumeSummary{Менторство}
          {менторил 3 стажёров, учил разработчиков проводить собесеодвания}
          {все стажёры были рекомендованы в штат, 1 стажёр остался в моей команде и дорос до middle, разработчики, которых я менторил, успешно проводят собеседования}
        \resumeSummary{Собеседования}
          {провожу технические и алгоритмические собеседования на \tech{Go} и \tech{Python}, собеседую стажёров в нашу команду}
          {есть кандидаты, которых наняли в нашу команду}
        \resumeSummary{Работа с командой}
          {вел проект от старта до запуска, в проекте было задействовано 3 разработчика. В процессе взаимодействовал со смежными менеджерами, пользователями и заказчиками, составлял план для разработчиков}
          {запустил большой проект, который используется на всех новых автономных грузовиках и такси Яндекса}
      \resumeItemListEnd

      \resumeSubheading
      {Школа программирования NlogN}{Москва, Россия}
      {Python Разработчик}{Июнь 2022 -- Сентябрь 2023}
      \resumeItemListStart
        \resumeItemDetailed{Платформа для сдачи заданий}
          {Разработал платформу для сдачи заданий: для ученика и для проверяющего.}
          {Ученики отправляли решения задач в личные сообщения преподавателю. Контроль вёлся с помощью google таблиц.}
          {TG-бот для ученика и проверяющего, где можно посмотреть список невыполненных заданий и отправить решение задачи. Автоматически обновляемая таблица с результатми и рейтингом всех учеников.}
          {Для ботов использовался \tech{aiogram} (\tech{python}), для работы с google таблицами использовался \tech{google-api-python-client}.}
        \resumeItemDetailed{Корпоративный мессенджер}
          {Развернул и поддерживал \tech{RocketChat}.}
          {Не у всех детей есть младших классов есть телеграм. При большом количестве групп в TG-аккаунте преподавателя начинается неразбериха.}
          {Развернул \tech{RocketChat}, перевёл взаимодействие туда. Добавил генерацию логинов и паролей для новых учеников, перевел TG ботов в \tech{RocketChat}.}
          {Взаимодействие с API \tech{RocketChat} из \tech{Python} с помощью \tech{rocketchat\_API}.}
        \resumeItemDetailed{Obsevrability систем команды}
          {Настроил отправку метрик и алёрты.}
          {Тысячи устройств для записи данных, состояние работы каждого из них необходимо отслеживать и вовремя реагировать на проблемы.}
          {Дашборды, метрики и алёрты для отслеживания состояния работы устройств и сервисов на них.}
          {Все метрики отправляются в Яндексовый аналог \tech{Prometheus}, отрисовываются в \tech{DataLens}, алёрты настроены в Яндексовом аналоге \tech{Grafana}.}
        \resumeItemDetailed{Сервис для контроля данных на флоте автономных машин}
          {Расширил функциональность: добавил логику сохранения данных при смене локации для конкретного робота, улучшил UI для настройки конфигов.}
          {При смене локации робота-доставщика требовалось ручное вмешательство. Редактирование конфигов было неудобным, приходилось копировать конфиг в сторонние редакторы.}
          {Придумал и реализовал логику работы сервиса при смене локации робота-доставщика. Подключил плагин для корректного отображения \tech{yaml} конфигов. Запуск новых локаций стал независимым от сервиса, время редактирования конфигов сократилось с 10 минут до 1 минуты.}
          {Бэкенд управления написан на \tech{Python} \tech{Django}. Метаданные хранятся в \tech{PostgreSQL} с \tech{master/replica}-роутингом,
          аналитика событий в \tech{ClickHouse}.}
        \resumeSummary{Менторство}
          {при увольнении онбордил нового сотрудника, заменяющего меня}
          {новый сотрудник продолжил поддерживать существующие сервисы}
      \resumeItemListEnd

      \resumeSubheading
      {Яндекс Облако, Маркетплейс}{Москва, Россия}
      {Стажёр-разработчик}{Октябрь 2022 -- Апрель 2023}
      \resumeItemListStart
        \resumeItemDetailed{Публикация OS-продуктов}
          {Реализовал пайп сборки и выкладки образов, написал документацию по процессу.}
          {Публикация партнёрских образов делалась вручную и требовала помощи команды.}
          {Публикация продукта самостоятельно через единый пайп.}
          {\tech{Python}, \tech{Packer}, \tech{Salt}}
        \resumeItemDetailed{Расширение API сервиса}
          {Добавил валидацию запросов и новый тип поля, провёл миграцию БД с бэкфиллом исторических данных.}
          {Часть полей не валидировалась; для нового функционала нужен был столбец с пересчётом старых строк.}
          {Некорректные запросы отбрасываются до бизнес-логики; миграция заполнила все исторические строки.}
          {\tech{Python}, \tech{PostgreSQL}}
        \resumeItemDetailed{Эксплуатация и observability}
          {Завёл \tech{monrun}-проверку в \tech{Salt}-формулу и добавил команду в \tech{Go} CLI команды.}
          {На стендах не было автоматической проверки нужного состояния; типовой запрос делался руками.}
          {Проблемы ловятся автоматически; типовой запрос делается одним вызовом из терминала.}
          {\tech{Go}, \tech{monrun}, \tech{Salt}}
      \resumeItemListEnd

  \resumeSubHeadingListEnd


%-----------PROJECTS-----------------
\section{Выступления/Лекции}
  \resumeSubHeadingListStart
    \resumeSubItem{Центральный Университет}
      {Веду курс по продуктовой разработке}
    \resumeSubItem{Яндекс.Железо}
      {Рассказывал про эволюцию продуктов нашей команды: \href{https://youtu.be/yXUiDlOObgo}{\underline{youtube}}}
    \resumeSubItem{Школа Бэкенд Разработки Яндекса}
      {Прочитал лекцию про дебаг логгирование и профилирование на \tech{python}}
  \resumeSubHeadingListEnd  

%-----------PROJECTS-----------------
\section{Проекты}
  \resumeSubHeadingListStart
    \resumeSubItem{Outline Work}
      {Сервис для написания конспектов. Релизовал приложение для работников и администраторов.}
    \resumeSubItem{Стоматология}
      {Разработал админку для автоматической работы со страховыми компаниями.}
    \resumeSubItem{TG-бот для получения кадастровой информации}
      {Разработал бота для получения кадастровой информации из Росреестра.}
  \resumeSubHeadingListEnd

%-----------EDUCATION-----------------
\section{Образование}
  \resumeSubHeadingListStart
    \resumeSubheading
      {Университет МИСИС}{Москва, Россия}
      {ИВТ, Бакалавриат}{2022 -- 2026}
    \resumeSubheading
      {Тинькофф Поколение, Алгоритмы и структуры данных}{Москва, Россия}
      {Дивизион B}{}
    \resumeSubheading
      {Яндекс Лицей}{}
      {Python}{}
    \resumeSubheading
      {IT-школа Samsung}{}
      {Java, Android}{}
  \resumeSubHeadingListEnd

%-------------------------------------------
\end{document}
